\documentclass[10pt,a4paper]{amsart}
\usepackage[margin=19mm]{geometry}
\usepackage{amsmath,amssymb,mathtools}
\usepackage[scr=rsfs,cal=euler]{mathalfa}
\usepackage{xfrac}
\usepackage[unicode,hidelinks,bookmarks=false]{hyperref}
\newcommand{\ie}{i.e.\ }
\newcommand{\cf}{cf.\ }
\newcommand{\resp}{respectively\ }
\newcommand{\Subsection}{\subsection}
\providecommand{\nobreakdash}{\nobreak\hbox{-}}
\setlength{\emergencystretch}{2em}
\multlinegap0pt
\usepackage{amssymb}
\usepackage{mathtools}
\usepackage[shortlabels]{enumitem}
\usepackage{float}
\newtheorem{theorem}{Theorem}
\title{On the smoothness of 3-dimensional \\ skew polynomial rings}
\author{Andr\'es Rubiano, Armando Reyes}
\date{Source: arXiv:2603.11255v1 (CC BY 4.0)}
\begin{document}
\maketitle

\noindent\emph{Source and licence.} On the smoothness of 3-dimensional \\ skew polynomial rings. Version arXiv:2603.11255v1 (CC BY 4.0), distributed by arXiv under the Creative Commons Attribution 4.0 International licence, http://creativecommons.org/licenses/by/4.0/, as recorded in the arXiv licence metadata for this exact version and shown on the abstract page https://arxiv.org/abs/2603.11255v1. Full licence notice: \texttt{licenses/arxiv-2603.11255-CC-BY-4.0.txt}.\par\smallskip

\noindent\textit{Editorial scope.} Counted unit: the article's theorem Secondtheoremsmoothness3-dimensionalalgebras (source0.tex lines 990--992). The article's complete original proof of it is delivered with it (source0.tex lines 993--999, one \texttt{proof} environment of 7 lines). No mathematics is written, repaired or re-derived: the statement and the proof are the article's own lines, in the article's own notation, and no retained line is substituted or re-flowed.  Retained uncounted as the source-local context the proof needs: lines 564--568.  Nothing was omitted from any retained span, so the package carries no omission record.  No retained line contains a citation, so the wrapper carries no bibliography and no bibliography entry.  No retained line contains a figure environment, an image-inclusion command or drawing code, so no image is delivered and the package declares no asset dependency.  Editorial formatting only, in addition: an AMS \texttt{amsart} preamble that restates the article's own declarations and macros used by the retained lines, namely the environments \texttt{theorem} on separate counters, together with the standard packages those lines need.  The retained environments print in this document as Theorem 1 in order of appearance, whereas the article prints the same items with the numbering of its own full document.  The complete native source of the article as distributed in the arXiv e-print for this exact version is bundled unchanged as \texttt{sources/196204/source0.tex}.
\begin{align}
yz-\alpha zy & = a_{\lambda}x+b_{\lambda}y+c_{\lambda}z+d_{\lambda}, \notag \\
zx-\beta xz & = a_{\mu}x+b_{\mu}y+c_{\mu}z+d_{\mu}  \quad {\rm and}\ \label{rel3skew} \\
xy-\gamma yx & = a_{\nu}x+b_{\nu}y+c_{\nu}z+d_{\nu}, \notag
\end{align} 

\begin{theorem}\label{Secondtheoremsmoothness3-dimensionalalgebras}
If one of the conditions $a_{\lambda}\not = 0$, $b_{\mu}\not=0$ or $c_{\nu}\not=0$ holds, then there are no one-dimensional connected integrable calculi over $A$.
\end{theorem}

\begin{proof}
By contradiction. Suppose that $A$ has a first order differential calculus $(\Omega A, d)$. Without loss of generality, we consider the case $a_{\lambda}\not =0$. Since $d$ must be compatible with the relations {\rm (}\ref{rel3skew}{\rm )}, then we get that
\begin{equation*}
    dy z-\alpha ydz = a_{\lambda}dx +b_{\lambda}dy + c_{\lambda}dz,
\end{equation*}
whence $dx$ is generated by $dy$ and $dz$. This means that $\Omega^1A$ is generated by two elements and $\Omega^3A=\Omega^1A\wedge\Omega^1A\wedge\Omega^1A=0$, i.e. there is no third-order calculus. Since ${\rm GKdim}(A) = 3$, the assertion follows.
\end{proof}

\end{document}
